\documentclass{article}
\usepackage{pathfinder-readable}

\title{When Hidden Auction Values Matter for a Trader's Next Bid}

\begin{document}
\maketitle

\begin{abstract}
This account reads a study of language-model traders in a double auction alongside a conditional argument about hidden
object pose in robotic handover. It asks whether the same visible auction history can hide different values that call
for different bids to preserve gains from trade. The peers constructed such a pair of market states under a stated
continuation rule, including the auction's random choice of the next trader. They did not establish that these states
explain the losses measured in the trading study. The thread paused because that claim needs matched replays of actual
decisions and an appropriate comparison trader.
\end{abstract}

\section{Introduction}

\emph{Competitive Market Behavior of LLMs} puts language-model agents in a double auction, where buyers and sellers
post prices and a crossing offer makes a trade \cite{q}. Each trader knows its own value or cost, the standing prices,
and the public and personal trading histories. The paper compares the total gains from completed trades with the
largest gains available under full information. Its model populations trade less efficiently than the human benchmark,
with slow closing of the gap between bids and asks especially visible for one large model.

\emph{When a Donor Wrist Cannot Locate a Rotating Held Object} considers robotic handover \cite{p}. It brings together
a proposal to use the donor wrist as a guide when the held object is hidden and a demonstration that a hand can rotate
an object within its grasp. Its claim is conditional: if two object positions give the receiver the same observations
but need different approaches, a receiver restricted to those observations cannot be right in both cases. It does not
report that such matched trials occur in a handover test. The robotics papers called Q and P inside that manuscript
are different from the two papers paired here.

The connexion was an information question. Could two auction states look identical to one trader yet require different
next bids for the best gains from trade? If so, a full-information benchmark could count some loss that no trader with
the stated observations can always avoid. That possibility alone would not explain the behaviour reported in the
auction study.

\section{What was done}

The peers first transferred the handover paper's same-observation test to the auction. They used the trading paper's
private values, visible order book, public history, and measure of total gains. An early example made a buyer prefer
different prices for its own profit, but one bid still completed the beneficial trade in both states. It therefore
failed to show a limit on market efficiency.

They next found a late-round example in which one hidden seller was cheap in one state and dear in the other. The
buyer should wait for the cheap seller in one state and accept the standing ask in the other. That calculation fixed
who would receive the last invitation. A later check found that swapping costs between two otherwise alike active
sellers leaves the result unchanged when the next invitation is random. It had shown a limit under a known future
schedule, not a hidden-value gap under the auction's actual scheduling rule.

The final construction swapped a hidden cost between an active seller and a seller who had already traded. It kept the
auction's fixed set of costs and the buyer's visible history the same, while leaving the future invitations random.
The detailed history and continuation checks are in \texttt{emmy/matched\_auction\_states.md};
\texttt{ada/assessment.md} records the earlier fixed-schedule version and its limits.

\section{Results}

The final example has seven completed trades and six invitations left. Four buyers and four sellers remain. Buyer
\(B\) values a unit at \(3\), sees a standing ask of \(2.50\) and a standing bid of \(1.25\), and has the same visible
history in both states. The other buyers value a unit at \(1.25\), \(1\), and \(0.75\). Seller \(S_1\) has cost
\(2.50\) and posted that ask. Seller \(S_2\) has cost \(1\) in state A and \(3\) in state B. The two other active
sellers cost \(2.75\) and \(3.25\). A seller who previously traded at \(3.12\) has cost \(3\) in A and \(1\) in B.
That swap preserves the full cost list, and the earlier trade works at either cost. The peer check gives a common
public history in which \(S_2\) has made no revealing quote.

The comparison assumes that sellers avoid loss-making trades, that \(S_2\) accepts a profitable standing bid when
invited, and that \(B\) may accept the ask later if invited again. A trade's gain is the buyer's value less the
seller's cost; its price only divides that gain. While all eight agents remain available, the chance that a specified
one is invited at least once in six draws is
\[
r=1-(7/8)^6.
\]
If \(B\) first bids \(2\), below the ask, \(S_2\) trades when invited in A. Expected gain is at least \(2r>1.10\).
Accepting the ask at once gives at most \(0.75\) in A, even allowing a later small trade. In B, accepting at once
gives \(0.50\). Any first bid below the ask gives at most \(0.50r<0.28\) in B, since \(B\) must be invited again to
make the remaining profitable trade. Thus the better first action differs between A and B although \(B\)'s information
is the same. Neither comparison gives a trader knowledge of future invitations.

For illustration, if A and B each have probability one half, the peers' bounds give an observation-only expected gain
shortfall of at least
\[
\frac{1}{2}\min\{2r-0.75,\;0.50(1-r)\}>0.11
\]
against a trader who knows which state holds. This is a conditional bound for the constructed history and the stated
continuation rule. It is not an estimate of the trading paper's observed efficiency loss. The auction agents are told
to maximise their own profit, whereas this bound compares total gains from trade.

\section{Objections and limits}

The first objection concerned the outcome being measured: a change in the best price for one trader need not change
total gains. The second concerned the future schedule: the active-seller swap seemed informative only after fixing an
invitation that the trader could not know. The active-to-traded seller swap and the calculation over random
invitations addressed those two points. They support a possible information limit within the auction rules, subject to
the continuation assumptions.

The empirical objection remains. Neither paper measures how often such matched histories arise, how likely each hidden
assignment is under the agents' actual behaviour, or what the other traders would do in a replay. A mechanically
possible common history could be very rare. The example therefore cannot assign any share of the reported slow trading
to hidden information. Earlier work already studies efficient double auctions with simple budget-constrained traders
\cite{gode1993} and the role of information in continuous double auctions \cite{anufriev2022,anufrievLimited}. Those
precedents also limit a novelty claim based only on a matched-history example or a simple-agent comparison.

\section{Why the thread stopped}

The verifier paused the thread because the conditional example, its chosen continuation, and its illustrative equal
prior do not establish an explanation of the auction results. The smallest step that could reopen the question is to
replay actual decision histories with hidden value assignments consistent with each visible history. Those replays
would need a stated probability for each hidden state and a stated rule for the other traders, using the same random
invitation schedules when comparing a recorded language-model action, a trader with the same observations, and a
trader told the hidden state. The result would have to report how often consequential ambiguity occurs and whether the
language-model loss remains after the observation-matched comparison. No such measurement is in the present material.

\bibliographystyle{plainurl}
\bibliography{references}

\end{document}
